% !TEX TS-program = pdflatex
% David - Asentaja CV (suomi)
% Layout from hmemcpy/cv template (adapted from Professional CV by Alessandro Plasmati).

\documentclass[a4paper,11pt]{article}

% Shared style: works from roles/<role>/ in this repo, or with cv-style.sty
% uploaded next to this file (e.g. in Overleaf).
\IfFileExists{cv-style.sty}{\usepackage{cv-style}}{\usepackage{../../shared/cv-style}}

\cvlang{fi}
\cvrole{Asentaja}

\begin{document}

% --------------------TITLE-------------
\begin{cvheader}
\par{\centering
		{\Huge \textsc{David Adegbola}}\\[0.4em]
		{\Large Asentaja}
	\bigskip\par}

\hrule
\vspace{0.5em}
\begin{tabular}{rl}
  \textsc{\cvt{location}:} & Joensuu, Suomi (valmis muuttamaan Iisalmeen)\\
  \textsc{\cvt{phone}:}    & +358 46 800 7210\\
  \textsc{\cvt{email}:}    & \href{mailto:adavidadegbola@gmail.com}{adavidadegbola@gmail.com}\\
\end{tabular}
\end{cvheader}

\section{\cvt{summary}}
\begin{tabular}{p{0.9\textwidth}}
 Hei! Olen David, 24-vuotias kestävän metsäbiotalouden opiskelija Itä-Suomen yliopistossa. Muutin hiljattain Vaasasta Joensuuhun opintojeni vuoksi. Pidän käytännön työstä, jossa näen kätteni jäljen. UPM:n Kaukaan sellutehtaalla tein asennustöitä sähkösuodattimen uusimisessa: käytin kulmahiomakonetta ja nostolaitteita, pidin työmaan siistinä ja noudatin tehtaan tiukkoja turvallisuusohjeita. Työnantaja arvioi työni kiitettäväksi. Olen myös jakanut lehtiä yksin aikaisin aamulla ja pystyttänyt AV-laitteita joka viikonloppu. Näissä töissä olen oppinut tarkkuutta, huolellisuutta ja turvallisia työtapoja. Olen oma-aloitteinen ja toimin hyvin sekä yksin että tiimissä.\\\\

 Haluan oppia sähkö-, hydrauliikka- ja mekaanista asennusta, ja olen valmis muuttamaan Iisalmeen ja tekemään kokoaikaista päivätyötä. Suomen kielen taitoni on vielä perustasolla, mutta opiskelen sitä aktiivisesti, ja asentajan työ olisi hyvä paikka harjoitella sitä käytännössä.
\end{tabular}

\renewcommand{\arraystretch}{1.12}
\section{\cvt{skills}}
\begin{tabular}{rl}
  \textsc{Asennus:} & Asennus- ja kokoonpanotyöt; kulmahiomakone, nostolaitteet, sähkötyökalut\\
  \textsc{Työtavat:} & Työohjeiden noudattaminen, työkalujen ja materiaalien inventaario, työturvallisuus\\
  \textsc{\cvt{cards}:}     & Työturvallisuuskortti\\
  \textsc{\cvt{languages}:} & Suomi (perustaso, opiskelen aktiivisesti), englanti (sujuva), yoruba (äidinkieli)\\
  \textsc{Tietotekniikka:} & Tietotekniikan perustaidot\\
\end{tabular}

% Dates are split over two rows so the long Finnish month names fit
\section{\cvt{experience}}
\begin{longtable}{r|p{0.72\textwidth}}

  \textsc{18.8.2025 --} & \job{Asentaja}{Axioma Servisas UAB}\\
  \textsc{7.10.2025}
    &UPM Kaukaan sellutehdas, Lappeenranta\\
    &Asennustyöt sähkösuodattimen uusimisessa.\\
    &Kulmahiomakoneen ja nostolaitteiden käyttö.\\
    &Työmaan siisteys sekä työkalujen ja materiaalien seuranta.\\
    &Tehtaan tiukkojen turvallisuusohjeiden noudattaminen.\\
    &Työnantajan arvio: kiitettävä. Vastuullinen, tunnollinen, halukas oppimaan ja ottamaan vastuuta.\sotagrow{asennus,kulmahiomakone,nostolaitteet,työturvallisuus,tarkkuus}\sep

  \textsc{4/2025--8/2025} & \textbf{Lehdenjakaja}, Finnbox ja SSM, Vaasa: jakelureitit yksin aikaisin aamulla.\\[0.4em]
  \textsc{11/2024--8/2026} & \textbf{AV-vapaaehtoinen}, Vaasa City Church: ääni- ja videolaitteet joka viikonloppu.\\[0.4em]
  \textsc{kesä 2026} & \textbf{Founder in Residence}, Vaasa Startup Factory: tuotekehitystä tiimissä.\\[0.4em]
  \textsc{2016--2020} & Myynti- ja asiakaspalvelutyötä Nigeriassa (Desire Feeds, Bunmi Pastries).\\
\end{longtable}

\section{\cvt{education}}
\begin{tabular}{rl}
  \textsc{2026--nyt:}& B.Sc, Kestävä metsäbiotalous, Itä-Suomen yliopisto\\
  \textsc{2024--2026:}& Tietotekniikan insinööriopintoja (keskeytetty), Vaasan ammattikorkeakoulu\\
  \textsc{2019--2022:}& OND, tietojenkäsittely, The Polytechnic, Ibadan, Nigeria\\
\end{tabular}

\end{document}
